% =====================================================================================
\chapter{RESULTS AND DISCUSSION}
\label{chp:results}
% Material: README §14 (the error chain and budget), §15-16; thesis_notes.md (every
% number with its date); scripts/foundationpose_results/{tack_marks.json, multiview/};
% todo.md "Benchmarks to run" and "measure everything" (sections marked PLANNED).
% =====================================================================================

\TODO{The setup and the procedure are Chapter~\ref{chp:setup}.}

\section{Calibration Results}
\label{sec:exp_calib}
\begin{table}[htbp]
\caption{Camera depth against the pen-measured table, three ranges (mm).}
\label{tab:exp_depth}
\centering
\begin{tabular}{@{}lccc@{}}
\toprule
 & near & middle & far \\
\midrule
Before calibration & +1.92 (291) & +0.31 (439) & $-$3.28 (626) \\
Tare, distance entered by hand & +1.71 (285) & +3.25 (423) & +5.24 (604) \\
Tare, distance computed from the robot & +0.29 (297) & +1.06 (443) & +1.43 (657) \\
\bottomrule
\end{tabular}
\par\smallskip\footnotesize Offset of the camera's table from the pen-measured plane in mm (range in mm), one series per state, 2026-10-02.
\end{table}
\TODO{Tilt separation (camera-fixed 0.49$^\circ$ / table 0.44$^\circ$ $\to$ 0.07$^\circ$);
hand-eye residuals; TCP check.}
\SRC{README \S14, \S16; thesis\_notes.md.}

\section{Registration Accuracy}
\label{sec:exp_registration}
\TODO{The normal gate (lean 6.5$^\circ\to$0.3--0.8$^\circ$, root 10 mm $\to\pm$0.2 mm);
pose jitter of the ICP tracker on a still part (2026-09-25: 5.2 / 2.5 / 3.0 mm std, 24 mm range,
swings to 6.4$^\circ$), which motivated the robust mean (FoundationPose's jitter on a still part is
not measured yet: planned with pose\_jitter\_probe on /perception/fp/pose); the save-time ICP
(synthetic: 4.7 mm / 1.5$^\circ\to$ 0.01 mm along the normal and 0.4$^\circ$ of tilt, the remaining
1--2 mm in the plate's own plane; bench: corrections 1.8--7.9 mm).}

\section{Multi-View Refinement}
\label{sec:exp_multiview}
\TODO{Synthetic: recovery of an injected extrinsic error (to 0.2 mm), acceptance logic.
Real: corrections per part, the estimated $\vect{d}$ across days, fit-up gaps, the
table offset per view before / after the correction (6.9 $\to$ 3.7 mm spread), the
prior sweep (3 / 8 / 15 mm: $<$0.5 mm change).}
\begin{table}[htbp]
\caption{Table height seen by each view relative to the pen-measured plane (mm).}
\label{tab:exp_table_views}
\centering
\begin{tabular}{@{}lccccc@{}}
\toprule
 & view 0 & view 1 & view 2 & view 3 & spread \\
\midrule
Raw & +2.46 & $-$1.62 & +2.64 & +5.26 & 6.9 \\
$\mathbf{R}_v\vect{d}$ taken out & $-$1.07 & $-$3.89 & $-$1.63 & $-$0.21 & 3.7 \\
\bottomrule
\end{tabular}
\end{table}
\SRC{capture multiview/20261006-163931; multiview\_refine\_offline.py; todo.md.}

\section{Placement Accuracy}
\label{sec:exp_placement}
\begin{table}[htbp]
\caption{The error chain: each cause found, how it was measured, and the mark error after
its fix.}
\label{tab:exp_error_chain}
\centering\small
\begin{tabular}{@{}p{0.09\textwidth}p{0.33\textwidth}p{0.33\textwidth}p{0.15\textwidth}@{}}
\toprule
Date & Cause & Fix & Mark error after \\
\midrule
24 Sep & failed hand-eye solve; TCP frame & Park--Martin re-solve; TCP composed & $\sim$8 mm \\
25 Sep & tracker jitter at save & robust mean of poses at rest & 5--8 mm \\
28 Sep & flat ground cut on a tilted table & measured table plane & 3--5 mm \\
28 Sep & camera-fixed tilt 0.36$^\circ$ & tilt refinement & 3.7--4.3 mm \\
29 Sep & nominal kinematics; plate face straddle & calibrated chain; normal gate & \\
1--2 Oct & depth $\propto z^2$ (contact 8.3--8.8 mm early) & Tare against the robot & \\
2 Oct & extrinsic translation; single view & multi-view refinement & $\sim$2.5 mm \\
6 Oct & tracker pose saved (FoundationPose) & ICP at save; APS; roll fallback & 2.3--2.7 mm \\
\bottomrule
\end{tabular}
\end{table}
\SRC{README \S14 tables; thesis\_notes.md.}
\TODO{Rows missing against README \S14 / thesis\_notes.md (review B2, @X ch6): 29 Sep the holder
envelope (clearance 1.1 $\to$ 15.5 mm); 29 Sep an unmodelled 7 mm plate under the ear (9.7 mm early
contact); 2 Oct the extrinsic tilt 0.49$^\circ\to$0.07$^\circ$; and the empty ``after'' cells of
29 Sep (straddle: marks 14--16 mm before the fix) and 1--2 Oct. Rebuild the table from the dated
entries.}
\TODO{The error budget per link (README \S14), and repeated trials with statistics
(mean, std, max over $N$ tacks), not single runs.}

\section{Live Tracking: FoundationPose versus ICP}
\label{sec:exp_tracking}
\TODO{PLANNED (todo.md, Benchmarks): pose\_jitter\_probe on both trackers, still part
and slow push: noise, delay, drift; rate 25--29 Hz vs.\ $\sim$10 Hz.}

\section{Transit Planning: APS versus RRT-Connect}
\label{sec:exp_planning}
\TODO{PLANNED (aps\_transit\_plan.md step 5): per transit, planning time, weighted joint
travel, tool-tip path length, deviation, clearance, TOTG duration. Known so far:
front-back transit on the T-joint with $\sim$6$\times$ less weighted joint travel at a 1 s budget
(2.9 against 17.6--18.0, the median of five RRT-Connect runs); C++ collision check $\sim$590$\times$
faster (8.6~$\mu$s against 5.1~ms) and identical on 8\,800 configurations (2 $\times$ 4\,400, random
and near-contact, in two scenes).}
\figplaceholder{fig:exp_transits}{Tool-tip paths of RRT-Connect and APS for the
front--back transit (RViz or matplotlib, same view)}{Transit paths: RRT-Connect vs.\ APS.}

\section{Timing and Resources}
\label{sec:exp_timing}
\TODO{PLANNED (todo.md, ``measure everything''): wall and CPU time, memory, per stage
from the received pose to the mark; scaling with points, parts, views and transits.}

\section{Curved Parts}
\label{sec:exp_curved}
\TODO{The printed parts (C1, E2, S3 on R2, RR1, SP3) through the cell, single view and
multi-view (bench 2026-10-09..16: curved\_seams\_plan.md tables; thesis\_notes.md). Mark
errors per part and per view mode; the cut pipe's spin settled; the cable window.}

\section{Operator Interaction: the Camera Aim}
\label{sec:exp_aim}
\TODO{Measure after deciding the Jacobian fix (\S\ref{sec:mc_aim}): the achieved rate is half the
commanded one until kinematics.ur5e\_jacobian is fixed; the node logs the commanded rate.}
\TODO{Time-to-aim, overshoot, the stop when the hand drops (to measure). The command
staircase: the two bags (50 Hz direct vs.\ the 500 Hz relay) - commands, step size, the
wrist effort's 20--250 Hz ripple 3.4$\to$1.7$\times$ idle, the 50.8 Hz line gone
(scripts/aim\_bag\_report.py; thesis\_notes.md 2026-10-16). Figures: the effort spectrum
before/after; an 80 ms zoom of commanded vs.\ measured velocity.}

\section{WeldSet Benchmark}
\label{sec:exp_weldset}
\TODO{The seven literature methods and the two proposed extractors
(\S\ref{sec:ws_own}) per stratum and condition: F1, lateral p95, Chamfer; the oracle
ladder's L0$\to$L1 delta.}
\begin{table}[htbp]
\caption{Seam extraction per method and stratum (F1 at 3 mm / lateral p95, mm).}
\label{tab:ws_benchmark}
\centering\small
\begin{tabular}{@{}lccccc@{}}
\toprule
Method & T/line & corner & butt & lap & \ldots \\
\midrule
\textit{TBD} & & & & & \\
\bottomrule
\end{tabular}
\end{table}
\SRC{weld\_generator Phase 4 batch (desktop, after the fidelity audit); notebooks 15--17.}

\section{Discussion}
\label{sec:exp_discussion}
\TODO{What limits accuracy now (the remaining vertical offset, the pen length, unfixed
parts); what the measurements say about vision-only placement.}
